We first star briefly introducing a useful tool in the development of te formula we wish to use: the Bernoulli numbers $B_n$, polynomials $B_n(x)$ and periodic polynomials $\tilde{B}_n(x)$. The Bernoulli numbers $B_n$ are a sequence of rational numbers that frequently appears in the Taylor series expansion of trigonometric-type functions. These number can also be seen as special values obtained by the Bernoulli polynomials $B_n(x)$. This is because these polynomials can be written explicitly as
\begin{equation}
	B_k(x) = \sum_{n=0}^{k} {k \choose n} B_n x^{k-n},
\end{equation}
then, we must have that $B_k(0) = B_k$. More interestingly for us is the fact that these polynomials satisfy the differential equation 
\begin{equation}
	B'_n(x) = n B_{n-1}(x), \quad n = 1,2,\cdots
	\label{Equation: Differential Bernoulli}
\end{equation}
Finally, the periodic Bernoulli polynomial is simply the Bernoulli polynomials evaluated at the fractional part of the argument, i.e., $\tilde{B}_n(x) = B_n(x - \floor{x})$. Now, let $B_n$ and $B_n(x)$ denote the $n$th Bernoulli number and polynomial, respectively, and $\tilde{B}_n(x)$ the $n$th Bernoulli periodic function $B_n(x - \floor{x})$. Assume that $a, m$ and $n$ are integers such that $n > a$, $m > 0$, and $f^{(2m)}(x)$ is absolutely integrable over $[a,n]$. Then \cite{NISTMathFunc}, Let's start by

\begin{equation}
	\sum_{j=a}^{n} f(j) = \int_{a}^{n} f(x) \dd x + \frac{1}{2} \left(f(a) + f(n)\right) + \sum_{s=1}^{m-1} \frac{B_{2s}}{(2s)!} \left( f^{(2s-1)}(n) - f^{(2s-1)}(a)\right) + \int_{a}^{n} \frac{B_{2m} - \tilde{B}_{2m}(x)}{(2m)!} f^{(2m)}(x) \dd x.
\end{equation}

This formula is called the \textit{Euler-Maclaurin} formula. This equation can be extremely useful when estimating the asymptotic of the summation on the function $f$ when its integral can be more easily evaluate. In essence, this is a quadrature technique; a numerical integration: the approximation of an integral $\int f(x) \dd x$ of some function $f$ by a discrete summation $\sum w_i f(x_i)$ over a set of points $x_i$ with some weight function $w_i$. For example, in the classical Riemann sums, one approximates the integral $\int f(x) \dd x$ over some interval $[a,b]$ with the limit of the summation $\sum f(x_i^*) \Delta x_i$, where $x_i^* \in [x_{i-1}, x_i]$ and $ \Delta x_i = x_i - x_{i-1}$. Different methods can be generated by choosing $x_i^*$ differently.

The trapezoidal rule, for example, is a Riemann sum where $x_i^*$ is chosen to be the mean point $x_i^* = (x_{i-1} + x_i)/2$. In this way, we would estimate the area using trapezoids, as the name suggests. Then, one could write
\begin{align}
	\int_{a}^{n} f(x) \dd x & \approx \sum_{j = 0}^{n} f(x_j^*) = \sum_{j = 0}^{n} \frac{f(x_{i} + x_{i-1})}{2}, \\
	& = \sum_{j = 1}^{n-1} f(x_j^*) + \frac{1}{2} \left( f(n) + f(0) \right) = \sum_{j = 0}^{n-1} f(x_j^*) + \frac{1}{2} \left( f(n) - f(n) \right) .
\end{align}
Of course, this has some error term. Suppose that for some $(x_i, x_{i-1})$ we have the same $k$-derivatives at both extreme points for every $k\geq 1$, i.e., the function is constant. Then, there is no error in the approximation. If however, the two extreme points differ only on the $k$th derivative, there is some curvature defined by this difference where the error comes from. Therefore, there is indication of some error function given in respect to higher order derivatives.

Following the computations in \cite{Trapezoidal-Maclaurin}, take a function $f \in C^p(\R)$ for $p \geq 1$. To derive the Euler-Maclaurin for the interval $[a,b]$, let $h = (b - a)/n$ and $c = a + j h$, then we can express
\begin{equation}
	\int_{a}^{b} f(x) \dd x = \sum_{j=0}^{n-1} \int_{a+j h}^{a+(j+1)h} f(x) \dd x
\end{equation}
where we can use the property \eqref{Equation: Differential Bernoulli} and integration by parts to rewrite each integral as
\begin{align}
	\int_{c}^{c+h} f(x) \dd x & = h \int_{0}^{1} B_0(1-x) f(c+xh) \dd x \\
	& = - h \frac{B_1(1-x)}{1!} f(c + xh) |_0^1 + h^2 \int_{0}^{1} \frac{B_1(1-x)}{1!} f'(c+xh) \dd x, \\
	& \vdots \\
	& = - \sum_{r = 0}^{p-1} \frac{h^{r+1} B_{r+1}(1-x)}{(r+1)!} f^{(r)}(c+xh) |_0^1 + h^{p+1} \int_0^1 \frac{h^{r+1} B_{p}(1-x)}{(p)!} f^{(p)}(c+xh) \dd x, \\
	& = h \frac{f(c) + f(c+h)}{2} - \sum_{r = 0}^{p-1} \frac{h^{r+1} B_{r+1}(1-x)}{(r+1)!} \left[ f^{(r)}(c+h) - f^{(r)}(c) \right] + h^{p+1} \int_0^1 \frac{B_{p}(1-x)}{(p)!} f^{(p)}(c+xh) \dd x. 
\end{align}
where one can start to notice the trapezoidal rule creeping out on the first term. Notice that it was necessary to use that $ - B_1(0) = B_1(1) = 1/2$ and $B_n(0) = B_n(1) = B_n$ for $n\neq1$. Returning to the summation, we can rearrange terms as
\begin{align}
	h & \left[ \frac{f(a)}{2} + f(a+h) + \cdots + f(b-h) + \frac{f(b)}{2} \right] = \frac{1}{2} (f(b) - f(a)) + \sum_{j=0}^{n-1} f(a + j h), \\
	& = \int_{a}^{b} f(x) \dd x + \sum_{r = 0}^{p-1} \frac{h^{r+1} B_{r+1}(1-x)}{(r+1)!} \left[ f^{(r)}(b) - f^{(r)}(a) \right] - h^{p+1} \int_0^1 \frac{B_{p}(1-x)}{(p)!} f^{(p)}(c+xh) \dd x. 
\end{align}
This is the Euler-Maclaurin formula, somehow an extension of the trapezoidal rule including the higher order terms.